%-------------------------------------------------------------------------
% RESEARCH REPORT — Why resumes get rejected and what gets them accepted
% Compile with pdfLaTeX (Overleaf works out of the box).
% Markdown twin: 05_REJECTION_RESEARCH.md
%-------------------------------------------------------------------------

\documentclass[11pt,a4paper]{article}

\usepackage[margin=1in]{geometry}
\usepackage[T1]{fontenc}
\usepackage[utf8]{inputenc}
\usepackage[english]{babel}
\usepackage{enumitem}
\usepackage{array}
\usepackage{longtable}
\usepackage{titlesec}
\usepackage{xcolor}
\usepackage[hidelinks]{hyperref}

\definecolor{accent}{HTML}{0E5484}
\definecolor{warn}{HTML}{A03000}
\definecolor{rulegrey}{HTML}{999999}

\titleformat{\section}{\Large\bfseries\color{accent}}{\thesection.}{0.5em}{}
\titleformat{\subsection}{\large\bfseries}{\thesubsection}{0.5em}{}

\setlist[itemize]{leftmargin=1.4em, itemsep=2pt, topsep=3pt}
\setlist[enumerate]{leftmargin=1.6em, itemsep=2pt, topsep=3pt}

\newcommand{\keyfinding}[1]{\par\vspace{3pt}\noindent\textcolor{warn}{\textbf{#1}}\par\vspace{2pt}}

\title{\textbf{Why Resumes Get Rejected --- and What Gets Them Accepted}\\[4pt]
\large A research report for Sachin Kumar's iOS/macOS job search}
\author{Compiled from 2026 recruiter surveys, ATS scan studies,\\
eye-tracking research, and 26 real iOS job descriptions}
\date{August 2026}

\begin{document}
\maketitle

\noindent\rule{\textwidth}{0.4pt}
\vspace{4pt}

\noindent\textbf{Scope.} This is the second research pass, focused specifically on the
\emph{reject versus accept} decision. Every claim is sourced. Where a finding
\textbf{contradicts} earlier advice given in \texttt{01\_ATS\_PLAYBOOK.md}, it is called out
explicitly --- two findings did.

\vspace{4pt}
\noindent\rule{\textwidth}{0.4pt}

\tableofcontents
\vspace{8pt}

%=========================================================
\section{The numbers on rejection}

\subsection{The dominant cause is not formatting}

\begin{center}
\begin{tabular}{p{9.2cm} r}
\hline
\textbf{Finding} & \textbf{Number} \\
\hline
Rejections caused by experience not matching the posting & \textbf{73\%} \\
Recruiters rejecting for impersonal / untailored applications & \textbf{84\%} \\
Recruiters rejecting for lack of customization specifically & 54\% \\
Rejected resumes with \textless 50\% of the posting's required keywords, \emph{even when the candidate had the matching experience} & \textbf{82\%} \\
Resumes scoring below the 75-point ``qualified'' threshold & 71.4\% \\
Rejected resumes with a critical format failure & 34\% \\
Silent-rejection rate in \textbf{software engineering} & \textbf{53.3\%} \\
\hline
\end{tabular}
\end{center}

\keyfinding{Read row four again.}
82\% of rejected resumes were missing more than half of the posting's keywords
\emph{despite the candidate actually having that experience}. That is a \textbf{translation
failure}, not a qualification failure --- and it is precisely the failure mode of the original
resume: real unit-testing and REST-API work existed, but neither appeared as a bullet.

\subsection{Volume is the silent killer}
\begin{itemize}
  \item A typical posting draws \textbf{250 applications}; engineering roles report
        \textbf{2{,}000+ in the first week}.
  \item Roughly \textbf{5 interviews result per 180 applicants}.
  \item \textbf{53.3\%} of software-engineering applicants are \emph{silently} rejected ---
        no response at all.
\end{itemize}

\keyfinding{You are not being rejected. You are being out-ranked and never reached.}
Different problem, different fix.

\subsection{The Google VP's three red flags}

Brian Ong, VP of Recruiting at Google (Forbes, April 2026):

\begin{enumerate}
  \item \textbf{``Always needing to be right'' / closed-mindedness.} On a resume this reads as
        stagnant progression and \emph{no certifications within the last three years} --- a
        signal that you stopped growing.

        \keyfinding{This applies to you directly.}
        Both of your certifications are from \textbf{2022 --- four years old}. To a
        Google-style reviewer that reads as ``stopped learning in 2022.'' Fix in Section 4.

  \item \textbf{Inability to handle ambiguity.} \emph{``The ability to thrive in ambiguity is
        probably the most critical lens.''} Red flag: no examples of pivoting, adapting, or
        operating without enough resources or information.

        You \emph{have} this --- the legacy-library migration, the MVVM to VIPER refactor, and
        building a diagnostics tool nobody asked for. It was simply never framed that way.

  \item \textbf{Waiting for leadership instead of being the leader.} \emph{``Deferring to
        hierarchy all the time.''} On the page this shows up as passive verbs:
        \textbf{``supported,'' ``assisted,'' ``contributed to,'' ``worked on,'' ``helped
        with.''}

        \keyfinding{Your original resume opened a bullet with ``Contributed to architectural refactoring\ldots''}
        That is the exact passive construction named here.
\end{enumerate}

\subsection{Other documented rejection triggers}
\begin{itemize}
  \item \textbf{Unprofessional email} --- 76\% of such resumes ignored.
        \emph{You are fine.}
  \item \textbf{Photo, date of birth, marital status, father's name} --- legacy Indian biodata
        fields. They break parsing, add zero value, create bias risk, and some international
        ATS platforms \textbf{flag resumes containing photos as non-compliant}.
        \emph{You have none of these. Keep it that way.}
  \item \textbf{No side projects} --- 82\% of resumes have none. \emph{You have three, one
        live on the App Store. A genuine top-18\% advantage that was buried.}
  \item \textbf{Job hopping / no tenure} --- \emph{You are fine: one employer, continuous,
        with a promotion.}
\end{itemize}

%=========================================================
\section{What actually gets a YES}

\subsection{The 7.4-second scan, measured}

The Ladders eye-tracking study put recruiters in front of real resumes with an eye tracker.
Average initial screen: \textbf{7.4 seconds} (up from 6.0 in 2012). The measured gaze path:

\begin{enumerate}
  \item \textbf{Current job title and company} --- \emph{most fixation time of any element}
  \item Previous title and company
  \item Right side: employment \textbf{dates}, checking for steady progression
  \item Down to \textbf{education}
  \item Then, only if it survived, actual bullet content
\end{enumerate}

Winning resumes had simple layouts, clear headings, bold titles, bulleted achievements and
generous white space. Losing resumes had clutter, no white space, \textbf{multiple columns}
and long sentences.

\keyfinding{Implication.}
Your job title is the single most-looked-at element on the page --- and the original resume
rendered it in \emph{italic small text on the second line} of the entry. Every variant now
places a bold title headline directly under your name.

\subsection{The top third decides everything}
\begin{itemize}
  \item Recruiters spend \textbf{80\% of their attention on the top third} of the page.
  \item Resumes with \textbf{measurable results in the top third are 40\% more likely to be
        shortlisted}.
\end{itemize}

\keyfinding{Implication.}
A summary that is merely a list of framework names wastes the most valuable real estate you
own. Numbers belong above the fold --- which is why several new variants open with a
\textbf{Key Achievements} block: three quantified lines before the work history begins.

\subsection{Tailoring, measured three separate ways}

\begin{center}
\begin{tabular}{l r r}
\hline
\textbf{Measure} & \textbf{Untailored} & \textbf{Tailored} \\
\hline
Interview rate & 2.07\% & \textbf{4.23\%} \\
Callbacks & 1$\times$ & \textbf{3$\times$} \\
Hiring managers who actively prioritise evidence of tailoring & --- & \textbf{55\%} \\
\hline
\end{tabular}
\end{center}

\subsection{Referrals dwarf everything else in this report}

\begin{center}
\begin{tabular}{l r}
\hline
\textbf{Path} & \textbf{Hire rate} \\
\hline
Job-board application & 7\% \\
Referral & \textbf{30\%} \\
\hline
\end{tabular}
\end{center}

Referrals are \textbf{7\% of applications but 30--50\% of all hires}. Referred candidates are
approximately \textbf{15$\times$ more likely} to be hired.

\keyfinding{This is larger than every resume optimisation in this folder combined.}
Eleven resumes cannot beat one warm introduction. Use the resumes to \emph{convert} referrals,
not to replace them.

%=========================================================
\section{Two corrections to earlier advice}

\subsection{Correction 1 --- Two pages now beat one page}

The \texttt{OnePage} variant was built on the older ``one page under 8 years'' convention.
The 2026 data reverses this:

\begin{center}
\begin{tabular}{p{10cm} r}
\hline
\textbf{Finding} & \textbf{Number} \\
\hline
Recruiters who now \textbf{prefer two-page} resumes & \textbf{68.6\%} \\
Recruiters who still prefer one page & 21.6\% \\
HR professionals saying 1--2 pages is ideal ($n$ = 1{,}013) & 82.1\%, with \textbf{51\% preferring two} \\
Relative preference for two-page over one-page & \textbf{2.3$\times$} \\
Two-page resumes rated higher on ``comprehensive summary of the candidate'' & \textbf{+21\%} \\
Time recruiters spend reading a two-page resume & \textbf{more than 2$\times$ longer} \\
\hline
\end{tabular}
\end{center}

Mid-career (3+ years) is explicitly called out as needing two pages, and \textbf{98.8\% of
recruiters value clean formatting} --- they would rather you use a second page than cram
page one.

\keyfinding{What changes.}
\texttt{Sachin\_Resume\_OnePage.tex} is now a \emph{situational} tool: referrals, career
fairs, postings that explicitly say ``1 page,'' and US early-career applications. \textbf{It
is not your default.} Your default is a two-page variant.

\subsection{Correction 2 --- Certification recency is a scored signal}

The 2022 Coding Ninjas certificates were previously treated as neutral filler. Per Google's
VP of Recruiting, \emph{no recent certifications within three years} is read as evidence that
a candidate stopped growing. Yours are four years old. This is a real, fixable negative.

%=========================================================
\section{Your specific action list}

Ranked by impact per hour spent.

\begin{enumerate}
  \item \textbf{Get referrals.} 30\% versus 7\% hire rate. For every target company, find one
        42Gears alum or second-degree LinkedIn connection. This outweighs everything else here.
  \item \textbf{Kill every passive verb.} Google's VP names ``supported'' and ``assisted'' as
        an instant red flag; your original said \emph{``Contributed to architectural
        refactoring.''}
  \item \textbf{Earn one recent credential --- this year.} Cheapest honest options, in order:
        \begin{itemize}
          \item Ship the \textbf{Swift 6 strict-concurrency migration on CodeForge} and write
                a short public post about it. This simultaneously closes your \#2 skills gap,
                proves recent learning, and demonstrates thriving in ambiguity. Best single
                move on this list.
          \item An Apple Developer / Swift certification, or a reputable Swift Concurrency
                course with a dated certificate.
          \item A conference talk, blog series, or open-source contribution dated 2026 ---
                demonstrated growth counts, not only formal certificates.
        \end{itemize}
  \item \textbf{Put numbers in the top third.} Use the \texttt{Key Achievements} variants.
  \item \textbf{Make your job title impossible to miss.} Most-fixated element in the
        eye-tracking study.
  \item \textbf{Frame the ambiguity stories.} You have three. Do not just list them --- say
        what was unclear and what you decided.
  \item \textbf{Apply within 72 hours of posting.} With 2{,}000 applications in week one,
        queue position is a real variable.
  \item \textbf{Never send the same PDF twice.} 84\% of recruiters reject impersonal
        applications.
\end{enumerate}

%=========================================================
\section{India-specific findings}

\subsection{Naukri and Indian portals}
\begin{itemize}
  \item Naukri runs \textbf{its own parser and search ranking} before any company ATS sees
        you. Tables, columns, text boxes and graphics scramble it.
  \item Use \textbf{exact phrases} recruiters type: \texttt{SwiftUI} not ``Swift UI'';
        \texttt{Objective-C} not ``Obj C''; \texttt{CI/CD} not ``continuous integration''.
  \item Profile order that ranks: headline $\rightarrow$ current role $\rightarrow$ key skills
        $\rightarrow$ summary $\rightarrow$ resume upload.
  \item Razorpay, Swiggy and Zomato now receive thousands of applications per role through
        Naukri-integrated AI screening.
\end{itemize}

\subsection{Service companies versus product companies}
\begin{itemize}
  \item \textbf{IT services} (TCS, Infosys, Wipro, Cognizant, Capgemini, Accenture,
        LTIMindtree) hire on a bench-to-project model. They want \textbf{each project listed
        with stack, your role, team size, and client industry}. A senior Big-4-IT HR quote:
        \emph{``I can reject 80\% of resumes in 30 seconds because they're missing one of
        three things: 10th--12th marks, project details, or backlog status.''}
        Infosys uses \textbf{iRecruit + SmartRecruiters}; Wipro uses \textbf{Phenom People}.
  \item \textbf{Product companies} want impact, scale, ownership and metrics --- the format
        you already have.
\end{itemize}

\subsection{Applying abroad from India}
\begin{itemize}
  \item Add \textbf{one line of context} for names a foreign recruiter will not recognise.
        Recruiters abroad do not know 42Gears or CGC Landran.
  \item \textbf{US:} achievement-focused, 1--2 pages, ATS-critical.
  \item \textbf{Europe:} called a ``CV'', two pages standard, profile summary expected,
        academics weighted more heavily.
  \item State \textbf{work authorisation and timezone} explicitly for remote roles --- it is a
        knockout question, and knockout questions are what genuinely auto-reject.
\end{itemize}

%=========================================================
\section{The honest bottom line}

You asked for a resume that no recruiter can ignore. Here is what the data actually supports:

\begin{itemize}
  \item A great resume moves you from roughly \textbf{2\%} to \textbf{4--5\%} response on cold
        applications.
  \item A \textbf{referral} moves you to \textbf{30\%}.
  \item Nothing moves you to 100\%. With 2{,}000 applicants per role and a 53.3\%
        silent-rejection rate in software engineering, some strong candidates are never
        opened. That is arithmetic, not a flaw in your resume.
\end{itemize}

What this resume system \emph{does} guarantee: when a human or an LLM does open yours, there
is no reason to put it down. No parsing failure, no unevidenced claim, no passive verb, no
buried metric, no missing title. That is the entire controllable surface. Everything beyond
it is referrals, timing and volume.

%=========================================================
\section{Sources}

\subsection*{Rejection data}
\begin{itemize}\setlength\itemsep{1pt}
  \item \href{https://www.resumeadapter.com/blog/2026-ats-rejection-report-10000-resume-scans}{ResumeAdapter --- 2026 ATS Rejection Report, 10,000 resume scans}
  \item \href{https://resume.io/blog/resume-statistics}{Resume.io --- 58 resume statistics 2026}
  \item \href{https://www.jobsprout.ai/blog/resume-mistakes}{JobSprout --- 10 resume mistakes that get you rejected}
  \item \href{https://hirocv.com/blog/resume-statistics-2026-data-hiring}{HiroCV --- 50+ resume statistics for 2026}
  \item \href{https://enhancv.com/blog/ai-hiring-statistics/}{Enhancv --- AI hiring statistics 2026}
  \item \href{https://www.metaintro.com/blog/real-reasons-qualified-applicants-rejected-how-to-fix}{Metaintro --- Why qualified applicants still get rejected}
\end{itemize}

\subsection*{Recruiter behaviour}
\begin{itemize}\setlength\itemsep{1pt}
  \item \href{https://www.theladders.com/static/images/basicSite/pdfs/TheLadders-EyeTracking-StudyC2.pdf}{Ladders eye-tracking study (PDF)}
  \item \href{https://www.hrdive.com/news/eye-tracking-study-shows-recruiters-look-at-resumes-for-7-seconds/541582/}{HR Dive --- recruiters look at resumes for 7 seconds}
  \item \href{https://www.forbes.com/sites/rachelwells/2026/04/06/i-spoke-to-a-google-exec-he-reveals-the-3-resume-red-flags-that-instantly-get-you-rejected/}{Forbes --- Google exec on 3 resume red flags}
  \item \href{https://brainmanager.io/blog/career/resume-green-flags}{Brainmanager --- 20 resume green flags}
  \item \href{https://talent.seek.com.au/hiring-advice/article/green-flags-to-look-for-in-candidates}{SEEK --- green flags to look for in candidates}
\end{itemize}

\subsection*{Resume length}
\begin{itemize}\setlength\itemsep{1pt}
  \item \href{https://resumeoptimizerpro.com/blog/ideal-resume-length}{Resume Optimizer Pro --- ideal resume length}
  \item \href{https://www.gainrep.com/resources/one-page-resume-or-two/}{Gainrep --- one page or two, the 2026 recruiter verdict}
  \item \href{https://resumeworded.com/can-a-resume-be-two-pages-key-advice}{Resume Worded --- can a resume be two pages}
\end{itemize}

\subsection*{Referrals}
\begin{itemize}\setlength\itemsep{1pt}
  \item \href{https://www.zippia.com/advice/employee-referral-statistics/}{Zippia --- 25 employee referral statistics}
  \item \href{https://erinapp.com/blog/enterprise-employee-referral-statistics-you-need-to-know-for-2026/}{ERIN --- enterprise referral statistics 2026}
  \item \href{https://blog.theinterviewguys.com/referred-candidates-get-hired-4x-more-often-so-stop-asking-for-a/}{The Interview Guys --- referred candidates get hired 4x more often}
\end{itemize}

\subsection*{Big tech and startup}
\begin{itemize}\setlength\itemsep{1pt}
  \item \href{https://www.techinterviewhandbook.org/resume/}{Tech Interview Handbook --- FAANG-ready resumes}
  \item \href{https://mirrorcv.com/resume-guide/software-engineer-faang}{MirrorCV --- FAANG SWE resume guide 2026}
  \item \href{https://interviewkickstart.com/blogs/articles/software-engineer-resume-for-faang-companies}{Interview Kickstart --- SWE resume for FAANG}
  \item \href{https://www.recruitingfromscratch.com/blog/how-to-get-a-founding-engineer-job}{Recruiting From Scratch --- what startups screen for}
  \item \href{https://www.sweresume.app/articles/startup-resume-guide/}{SWE Resume --- startup resume guide}
\end{itemize}

\subsection*{India}
\begin{itemize}\setlength\itemsep{1pt}
  \item \href{https://cvefy.com/blog/ats-resume-format-naukri-india-guide}{CVefy --- ATS resume format for Naukri India}
  \item \href{https://www.cvforge.in/blog/resume-format-for-tcs-infosys-wipro}{CVForge --- resume format for TCS, Infosys, Wipro and Capgemini}
  \item \href{https://www.cv-prime.in/blog/ats-resume-mistakes}{CV Prime --- 15 ATS resume mistakes India 2026}
  \item \href{https://resumevera.com/guides/naukri-profile-optimization}{ResumeVera --- Naukri profile optimization 2026}
  \item \href{https://rezumea.com/resume-format/india}{Rezumea --- Indian resume format guide 2026}
  \item \href{https://cvcompose.com/us/blog/personal-information-resume-2026}{CVCompose --- personal information on a resume 2026}
\end{itemize}

\subsection*{International}
\begin{itemize}\setlength\itemsep{1pt}
  \item \href{https://resumeoptimizerpro.com/blog/understanding-international-resumes}{Resume Optimizer Pro --- international resume guide}
  \item \href{https://www.visualcv.com/blog/international-resume/}{VisualCV --- CV formats for 30+ countries}
  \item \href{https://www.infiniteresume.com/blog/indian-resume-international-jobs}{Infinite Resume --- Indian resume for international jobs}
\end{itemize}

\end{document}
